% Generated draft: requires mathematical and source review.
\documentclass[11pt]{article}
\usepackage[margin=25mm]{geometry}
\usepackage{fontspec}
\setmainfont{Latin Modern Roman}
\usepackage{amsmath,amssymb,mathtools,unicode-math}
\setmathfont{Latin Modern Math}
\usepackage{enumitem,needspace,xurl,hyperref}
\setlength{\parindent}{0pt}
\setlength{\parskip}{5pt}
\setlength{\emergencystretch}{3em}
\newcommand{\N}{\mathbb{N}}
\newcommand{\Z}{\mathbb{Z}}
\newcommand{\Q}{\mathbb{Q}}
\newcommand{\R}{\mathbb{R}}
\newcommand{\C}{\mathbb{C}}
\newcommand{\F}{\mathbb{F}}
\newcommand{\A}{\mathbb{A}}
\newcommand{\PP}{\mathbb{P}}
\newcommand{\E}{\mathbb{E}}
\newcommand{\eps}{\varepsilon}
\newcommand{\dd}{\,\mathrm{d}}
\begin{document}
\section*{UnsolvedMath ID: 701737}
\textbf{Source problem code:} catalogue entry
\subsection{Definitions and mathematical statement}
For a finite group \(K\), let \(\operatorname{Irr}(K)\) denote the set of irreducible complex characters of \(K\). For \(\chi\in\operatorname{Irr}(K)\), let \(\chi(1)\) be its degree and let \(\ker\chi\) be the kernel of a representation affording \(\chi\). Define the character codegree set by
\[
\operatorname{Cod}(K)=\left\{\frac{|K:\ker\chi|}{\chi(1)}:\chi\in\operatorname{Irr}(K)\right\}.
\]
This is a set, so repeated codegrees are counted only once. The supplied record does not define \(\operatorname{Cod}\); this definition uses its standard finite-group interpretation.

\textbf{Conjecture, under the finite-group interpretation.} For every finite non-abelian simple group \(G\) and every finite group \(H\),
\[
\operatorname{Cod}(G)=\operatorname{Cod}(H)\quad\Longrightarrow\quad G\cong H.
\]
Here a simple group has no normal subgroups other than the identity subgroup and itself.

\textbf{Unresolved domain qualification.} The supplied conjecture explicitly assumes that \(H\) is finite but calls \(G\) only a non-abelian simple group. Thus finiteness of \(G\) in the precise formulation above is an explicit interpretive hypothesis requiring confirmation. The supplied record gives neither a definition of \(\operatorname{Cod}(G)\) for infinite groups nor an explanation of whether infinite \(G\) are intended.

\textbf{Supplied comment of 2025.} The record reports that the conjecture has been proved when \(G\) is of type \(\operatorname{PSL}(2,q)\); \(\operatorname{PSL}(3,q)\) or \(\operatorname{PSU}(3,q)\); a sporadic simple group; an alternating group; \({}^{2}F_4(q^2)\); or \(\operatorname{PSL}(n,q)\) or a simple exceptional group of Lie type. The corresponding references are listed below. These are supplied dated literature claims, not independently verified conclusions. The record supplies no explicit parameter ranges or exceptional-parameter conventions for these families; its standing requirement that \(G\) be non-abelian simple must be retained.
\subsection{Short English statement}
G. Qian's conjecture asks whether a non-abelian simple group is determined, among finite groups, by its set of irreducible complex character codegrees. The finite-group formulation above explicitly assumes that the simple group is finite, a hypothesis left unstated in the supplied record. The record labels the problem open and includes a supplied 2025 comment reporting several proved families; neither the status nor those literature claims has been independently verified.
\subsection{Sources}
\begin{itemize}
\item Local UnsolvedMath registry: permanent ID \(701737\), exact source problem code \texttt{catalogue entry}, title ``Kourovka Notebook Problem 20.79.'' Supplied upstream revision: \texttt{d371ee917ac9647fbb06cd6fc529deec6f3dda82}. The record labels the problem \texttt{open} and records its update as \texttt{2026-10-16T14:00:00Z}; these are supplied registry metadata.
\item G. Qian, Problem 20.79, \emph{The Kourovka Notebook (Unsolved Problems in Group Theory)}, 20th issue (2022), with the supplied record also identifying the 21st issue edition (2026), \texttt{arXiv:1401.0300v48}. The collection description names E. I. Khukhro and V. D. Mazurov as editors. Attribution and edition information are supplied by the record; the source has not been independently consulted.
\item Supplied ``Comment of 2025'': the literature assessments below are supplied dated curation, not independently verified. No separate \texttt{literature\_checked\_at} date is supplied.
\item For \(\operatorname{PSL}(2,q)\): A. Bahri, Z. Akhlaghi, B. Khosravi, \emph{Bull. Austral. Math. Soc.} \textbf{104}, no. 2 (2021), 278--286.
\item For \(\operatorname{PSL}(3,q)\) or \(\operatorname{PSU}(3,q)\): Y. Liu, Y. Yang, \emph{Results Math.} \textbf{78}, no. 1 (2023), article No. 7.
\item For sporadic simple groups: M. Dolorfino, L. Martin, Z. Slonim, Y. Sun, Y. Yang, \emph{Bull. Austral. Math. Soc.} \textbf{109}, no. 1 (2024), 57--66.
\item For alternating groups: M. Dolorfino, L. Martin, Z. Slonim, Y. Sun, Y. Yang, \emph{Bull. Austral. Math. Soc.} \textbf{110}, no. 1 (2024), 115--120.
\item For \({}^{2}F_4(q^2)\): Y. Yang, \emph{J. Group Theory} \textbf{27}, no. 1 (2024), 141--155.
\item For \(\operatorname{PSL}(n,q)\) or simple exceptional groups of Lie type: H. P. Tong Viet, \emph{Math. Nachr.} \textbf{298}, no. 4 (2025), 1356--1369.
\end{itemize}
\end{document}
